\documentclass[10pt]{article}
\usepackage[left=1in,right=1in,top=3cm,bottom=0.35in,%
includeheadfoot,headheight=99pt,headsep=0.2in]{geometry}
\usepackage{times,graphicx,cite,fancyhdr,amsmath,amssymb,captdef}
\usepackage{marvosym}
\renewcommand{\headrulewidth}{0pt}
%----------------  page style -------------------
\pagestyle{fancy}
\fancyhead{}    %clear the head
\fancyhead[L]{
  \begin{minipage}[t]{1.8in}
    \hspace{0.6in}
    \includegraphics[height=0.6in]{anl.png}
  \end{minipage}
}
\fancyhead[C]{}
\fancyhead[R]{%
  \begin{minipage}[t]{1.7in}
    \begin{flushleft}
      {\scriptsize
        \vspace{-1.1in}
        {\bf Qingteng Zhang, Ph.D., Physicist}\\
        \vspace{0.08in}
        {\bf X-ray Science Division}\\
        Argonne National Laboratory \\
        9700 South Cass Avenue\\
        Argonne, IL 60439 \\
        \vspace{0.08in}
        \Telefon: 630-252-6366 \\
        }
    \end{flushleft}
  \end{minipage}
}
\fancyfoot{}    %clear the foot
\fancyfoot[C]{
  \footnotesize{A U.S. Department of Energy laboratory managed by UChicago
    Argonne, LLC}
}
% ==============================================================================
\setlength{\parskip}{2ex}
\setlength{\parindent}{0pt}

\usepackage[hidelinks]{hyperref}
\usepackage{enumitem}
\setlength{\parskip}{0.4em}
\setlength{\parindent}{0pt}
\raggedbottom
\sloppy   % long unbreakable e-mail addresses have no hyphenation points

\begin{document}

\today

Editors\\
\textit{Proceedings of the National Academy of Sciences}

Dear Editors:

Please consider the enclosed Research Report as a Direct Submission to
\textit{PNAS}: ``Dynamic Arrest Follows Assembly on a Temperature-Scaled
Common Clock in an Elastin-like Polypeptide.''

\textbf{Corresponding authors.} Qingteng Zhang
\mbox{(\href{mailto:qzhang234@anl.gov}{qzhang234@anl.gov})}, Arvind Ramanathan
\mbox{(\href{mailto:ramanathana@anl.gov}{ramanathana@anl.gov})}, and Gyorgy Babnigg
\mbox{(\href{mailto:gbabnigg@anl.gov}{gbabnigg@anl.gov})}.

\textbf{Classification.} Physical Sciences: Applied Physical Sciences.
Secondary: Biophysics and Computational Biology.

\textbf{Significance.} Many proteins assemble into dense phases that later
lose their mobility, but structure and mobility are usually measured with
different tools, so it has been hard to establish what sets the pace of arrest
and whether the growth of structure keeps pace with it.
We measure structure and motion on the same detector frames and find that they
run on one clock. In isothermal
experiments on the sequence-defined elastin-like polypeptide (VPAVG)$_{30}$ at
four temperatures between 28.00 and 28.75~$^\circ$C, every sample assembles,
but only those above 28.00~$^\circ$C arrest within the experiment. The
low-$Q$ SAXS intensity and the arrested fraction collapse with the same
scaling temperature. Along that clock, arrest follows assembly: in every sample
that arrests, the fast-decay amplitude at 175~nm falls to one half at nearly
the same low-$Q$ intensity, after most of the low-$Q$ structure has formed. Arrest reaches
the largest length scales first. At 175~nm the fast-decay amplitude falls from
0.80 to 0.07 within an hour, while the $1/e$ rate of the fast decay levels off.

\textbf{Novelty.} One rescaling of the waiting time collapses all four
isothermal experiments onto a single master curve: with $T_s=0.24\pm0.02$~K,
arrest speeds up by about a factor of three for every quarter of a degree of
heating, and the low-$Q$ SAXS superposes with the same $T_s$. That is a stronger statement than a rate: most of what
distinguishes the four experiments is their pace, and temperature sets it. That scaling temperature gives the clock an
activation enthalpy of 3.1~MJ/mol, 6 to 14 times that of the calorimetric
transition of the polymer, so arrest is set by an additional barrier rather than
by desolvation itself.

That scaling framework is our own. One of us introduced it for thermoreversible
colloidal gels [Q. Zhang \textit{et al.}, \textit{Phys. Rev. Lett.}
\textbf{119}, 178006 (2017)], where it returns $T_s=0.33$ and 0.45~K for two
particle sizes [D. Bahadur \textit{et al.}, \textit{J. Chem. Phys.}
\textbf{151}, 104902 (2019)]. The present work carries it from colloids to a
protein polymer and finds the same kind of steep exponential scaling. Time--temperature
superposition has also been followed with coherent X-rays in a heat-set protein
gel. In heated egg yolk, protein gelation and lipoprotein aggregation follow
Arrhenius clocks with activation energies of 410 to 450~kJ/mol, typical of
protein denaturation [N. D. Anthuparambil \textit{et al.}, \textit{Nat. Commun.}
\textbf{14}, 5580 (2023)]. We cite both the nanoparticle-gel phenomenology and
the egg-yolk superposition rather than claiming either as new. What is new is
the size of the clock's activation enthalpy, what shares the clock, and how it
compares with calorimetry. XPCS and calorimetry probe the same quantity, the
enthalpy difference between the transition state of a rate-limiting step and
its starting state. In a disordered protein with no folded structure to unfold,
the clock's activation enthalpy is 3.1~MJ/mol, about seven times that in egg
yolk and 6 to 14 times that of the calorimetric transition of poly(VPAVG)
[J. Reguera \textit{et al.}, \textit{Macromolecules} \textbf{36}, 8470 (2003)].
A value this large does not mean an impassable barrier, because a
matching entropy of activation offsets most of it. The arrest and the low-$Q$
structure return the same scaling temperature from the same detector frames,
while the fast relaxation at a given stage hardly changes with temperature, so
the activation enthalpy belongs to the evolution toward arrest. That evolution
is therefore set by an additional barrier rather than by desolvation itself. We
set out two readings: a barrier of another
origin, such as arrested phase separation, or a step that joins several chains.
In the second case the calorimetric enthalpy of one VPAVG repeat
[C.-H. Luan \textit{et al.}, \textit{Biopolymers} \textbf{31}, 465 (1991);
Reguera \textit{et al.}; A. Ribeiro \textit{et al.}, \textit{Biophys. J.}
\textbf{97}, 312 (2009)], interpolated to 30 repeats, counts about 330 to 440
pentapeptides, the content of 11 to 15 chains. Because the protein is
monodisperse and sequence-defined, isothermal series at several chain lengths,
with calorimetry of the same construct, could tell the two readings apart,
which a natural mixture such as egg yolk does not allow. For condensates, the
same measurement could tell whether a molecular nucleus or a collective step of
many contacts paces their maturation.

\textbf{Why it needed this measurement.} The fast relaxation takes from a
fraction of a millisecond to several milliseconds at these length scales, and
early in each experiment it is at or below the shortest delay a conventional
kilohertz acquisition reaches. Recording the scattering at 50~kHz with the
XSPA-500k at APS 8-ID-I resolves the fast decay itself: its rate and its
broadening as the sample arrests. At conventional speeds much of it would
appear only as a loss of contrast.

\textbf{Limits, stated plainly.} The 1.6~s maximum delay bounds a fraction
arrested on that scale rather than an infinite-time nonergodicity factor. The
scaling rests on three arresting temperatures spanning half a degree, and the
collapse is approximate rather than exact. The bond number applies only if
the rate-limiting step is built from the transition enthalpy that calorimetry
measures, the net of desolvation and chain contacts, and it rests on
calorimetry of longer poly(VPAVG) chains in water, interpolated to 30 repeats.

\textbf{Companion manuscript.} A separate manuscript, in preparation from the
same construct, reports the structural side: a thermally reversible transition
to a hierarchical assembly carrying $\beta$-sheet-like order, and the
dependence of that reversibility on chain length. The two were originally one
manuscript and were separated because the dynamical result stands on its own and
the structural result belongs with the chain-length series. There is no
overlap of figures or data between them. We are happy to supply the companion
manuscript on request.

\textbf{Data and code.} The reduced data and all analysis and figure-generation
code are public at
\href{https://github.com/qzhang234/PubCode-2021-IDPP}{qzhang234/PubCode-2021-IDPP};
every data-derived figure regenerates from that repository alone.

The manuscript is original, is not under consideration elsewhere and has been
approved by all authors. The authors declare no competing interest;
AI-tool use is disclosed in the Acknowledgments. Suggested Editorial Board
members, NAS members and reviewers are listed overleaf.

Thank you for your consideration.

Sincerely,

\vspace{0.1in}
\nopagebreak
Qingteng Zhang, Ph.D., on behalf of all authors

\newpage
\section*{Suggested Editors and Reviewers}
\small
PNAS asks for a minimum of three Editorial Board members, three NAS members and
three qualified reviewers. \textbf{TO BE COMPLETED BEFORE SUBMISSION.} The
researchers below have relevant expertise in elastin-like polypeptides,
protein and peptide assembly, small-angle X-ray scattering, and X-ray photon
correlation spectroscopy, and were carried over from the earlier cover letter;
they are reviewer suggestions, and the Editorial Board and NAS member lists
still have to be chosen.

\textbf{Bradley D. Olsen}\\
Alexander and I. Michael Kasser (1960) Professor of Chemical Engineering, Massachusetts Institute of Technology\\
\href{mailto:bdolsen@mit.edu}{bdolsen@mit.edu}\\
Expertise: elastin-like polypeptide gels, associative protein polymers, arrested phase separation, and neutron spin-echo measurement of relaxation in protein hydrogels.

\textbf{Darrin J. Pochan}\\
Professor of Materials Science and Engineering, University of Delaware\\
\href{mailto:pochan@udel.edu}{pochan@udel.edu}\\
Expertise: peptide and protein self-assembly, $\beta$-sheet hydrogel networks, and small-angle scattering and rheology of soft biomaterials.

\textbf{Andrei Fluerasu}\\
Lead Beamline Scientist, Coherent Hard X-ray Scattering, National Synchrotron Light Source II, Brookhaven National Laboratory\\
\href{mailto:fluerasu@bnl.gov}{fluerasu@bnl.gov}\\
Expertise: X-ray photon correlation spectroscopy and the microscopic dynamics of gels and arrested soft matter.

\textbf{Theyencheri Narayanan}\\
Senior Scientist, European Synchrotron Radiation Facility, Grenoble\\
\href{mailto:narayan@esrf.fr}{narayan@esrf.fr}\\
Expertise: small-angle X-ray scattering and X-ray photon correlation spectroscopy of colloids and soft matter.

\end{document}
